\section{Arquitectura de Control Propuesta}

La estrategia integra la estructura virtual, el consenso, la ley de control y la evasión de obstáculos en un esquema de control por capas.
\begin{figure}[H]
    \centering
    \begin{tikzpicture}[
      font=\fontfamily{ptm}\selectfont\scriptsize,
      >={Stealth[length=2mm]},
      block/.style={draw, minimum width=1.3cm, minimum height=1.3cm, align=center, inner sep=2pt},
    ]
      \node[block] (VS) {Estructura\\Virtual};
      \node[block, right=1.1cm of VS] (C)  {Consenso};
      \node[block, right=1.1cm of C]  (CL) {Ley de\\Control};
      \node[block, right=1.1cm of CL] (F)  {PVA};

      \node[align=center, above=0.7cm of C] (NB) {Vecindad};
      \node[align=center, above=0.7cm of F] (LI) {LiDAR};

      \draw[->] (VS) -- node[above]{$x_v$} node[below]{$[x,y]$} (C);
      \draw[->] (C)  -- node[above]{$z$} node[below]{$[a,\alpha]$} (CL);
      \draw[->] (CL) -- node[above]{$u_{ref}$} node[below]{$[v,\omega]$} (F);
      \draw[->] (NB) -- (C);
      \draw[->] (LI) -- (F);

      \coordinate (uEnd) at ($(F.east)+(0.8,0)$);
      \draw[->] (F.east) -- node[above]{$u^*$} node[below]{$[v,\omega]$} (uEnd);
    \end{tikzpicture}
    \caption{Estrategia de control de formación}
    \label{fig:flujo_control}
\end{figure}


%%%%%%%%%%%%%%%%%%%%%%%%%%%%%%%%%%%%%%%%%%%%%%%%%%%%%%%%%%%%%%

\section{Estructura Virtual}
\subsection{Generación de Referencias y Construcción del Error de Consenso}

La estructura virtual genera para cada agente la referencia asociada a su offset dentro de la formación.

\begin{center}
\begin{tikzpicture}[
  font=\fontfamily{ptm}\selectfont\scriptsize,
  >={Stealth[length=2mm]},
  block/.style={draw, minimum width=1.3cm, minimum height=1.3cm, align=center, inner sep=2pt},
]
  \node[block] (VS) {Estructura\\Virtual};
  \node[block, right=1.1cm of VS] (C)  {Consenso};
  \node[block, right=1.1cm of C]  (CL) {Ley de\\Control};
  \node[block, right=1.1cm of CL] (F)  {PVA};

  \node[align=center, above=0.7cm of C] (NB) {Vecindad};
  \node[align=center, above=0.7cm of F] (LI) {LiDAR};

  \draw[->] (VS) -- node[above]{$x_v$} node[below]{$[x,y]$} (C);
  \draw[->] (C)  -- node[above]{$z$} node[below]{$[a,\alpha]$} (CL);
  \draw[->] (CL) -- node[above]{$u_{ref}$} node[below]{$[v,\omega]$} (F);
  \draw[->] (NB) -- (C);
  \draw[->] (LI) -- (F);

  \coordinate (uEnd) at ($(F.east)+(0.8,0)$);
  \draw[->] (F.east) -- node[above]{$u^*$} node[below]{$[v,\omega]$} (uEnd);

  \draw[dashed, green!60!black] ($(VS.north west)+(-0.2,0.2)$) rectangle ($(VS.south east)+(0.2,-0.2)$);
  \draw[dashed, green!60!black] ($(C.north west)+(-0.2,0.2)$) rectangle ($(C.south east)+(0.2,-0.2)$);

  \coordinate (VSrectSouth) at ($(VS.south)+(0,-0.2)$);
  \coordinate (CrectSouth) at ($(C.south)+(0,-0.2)$);

  \draw[dashed, green!60!black, ->] (VSrectSouth) to[out=-60, in=240]
    node[below, pos=0.5] {\itshape offset del agente}
    (CrectSouth);
\end{tikzpicture}
\end{center}

\begin{center}
\footnotesize\textbf{Figura~\ref{fig:flujo_control}:} Del offset individual al error de consenso
\end{center}

%%%%%%%%%%%%%%%%%%%%%%%%%%%%%%%%%%%%%%%%%%%%%%%%%%%%%%%%%%%%%%

\section{Estructura Virtual}
\subsection{Generación de la Trayectoria del Grupo}

La trayectoria de la estructura virtual se genera mediante un spline de Catmull-Rom sobre un conjunto de waypoints $\{\mathbf{p}_k\}$. Para cuatro puntos de control consecutivos $(\mathbf{p}_0, \mathbf{p}_1, \mathbf{p}_2, \mathbf{p}_3)$, la posición interpolada para $t \in [0,1]$ es:
\begin{equation*}
    \mathbf{p}(t) = \tfrac{1}{2}\Big[2\mathbf{p}_1 + (-\mathbf{p}_0+\mathbf{p}_2)\,t + (2\mathbf{p}_0-5\mathbf{p}_1+4\mathbf{p}_2-\mathbf{p}_3)\,t^2 + (-\mathbf{p}_0+3\mathbf{p}_1-3\mathbf{p}_2+\mathbf{p}_3)\,t^3\Big].
\end{equation*}



\section{Estructura Virtual}
\subsection{Orientación de la Estructura Virtual}

La orientación $\theta_v$ se deriva de la tangente a la curva, calculada de forma analítica a partir del mismo polinomio:
\begin{equation*}
    \dot{\mathbf{p}}(t) = \tfrac{1}{2}\Big[(-\mathbf{p}_0+\mathbf{p}_2) + 2(2\mathbf{p}_0-5\mathbf{p}_1+4\mathbf{p}_2-\mathbf{p}_3)\,t + 3(-\mathbf{p}_0+3\mathbf{p}_1-3\mathbf{p}_2+\mathbf{p}_3)\,t^2\Big],
\end{equation*}
\begin{equation*}
    \theta_v = \operatorname{atan2}\big(\dot{y}(t),\, \dot{x}(t)\big).
\end{equation*}














%%%%%%%%%%%%%%%%%%%%%%%%%%%%%%%%%%%%%%%%%%%%%%%%%%
\section{Estructura Virtual}
\subsection{Coherencia entre la Referencia Generada y la Naturaleza del Movimiento}

Las referencias son consistentes con la naturaleza del movimiento del robot diferencial, evitando trayectorias que induzcan maniobras de reversa innecesarias.

\begin{center}
\begin{tikzpicture}[
  font=\fontfamily{ptm}\selectfont\scriptsize,
  >={Stealth[length=2mm]},
  block/.style={draw, minimum width=1.3cm, minimum height=1.3cm, align=center, inner sep=2pt},
]
  \node[block] (VS) {Estructura\\Virtual};
  \node[block, right=1.1cm of VS] (C)  {Consenso};
  \node[block, right=1.1cm of C]  (CL) {Ley de\\Control};
  \node[block, right=1.1cm of CL] (F)  {PVA};

  \node[align=center, above=0.7cm of C] (NB) {Vecindad};
  \node[align=center, above=0.7cm of F] (LI) {LiDAR};

  \draw[->] (VS) -- node[above]{$x_v$} node[below]{$[x,y]$} (C);
  \draw[->] (C)  -- node[above]{$z$} node[below]{$[a,\alpha]$} (CL);
  \draw[->] (CL) -- node[above]{$u_{ref}$} node[below]{$[v,\omega]$} (F);
  \draw[->] (NB) -- (C);
  \draw[->] (LI) -- (F);

  \coordinate (uEnd) at ($(F.east)+(0.8,0)$);
  \draw[->] (F.east) -- node[above]{$u^*$} node[below]{$[v,\omega]$} (uEnd);

  \draw[dashed, green!60!black] ($(VS.north west)+(-0.2,0.2)$) rectangle ($(VS.south east)+(0.2,-0.2)$);
  \draw[dashed, green!60!black] ($(CL.north west)+(-0.2,0.2)$) rectangle ($(CL.south east)+(0.2,-0.2)$);

  \draw[dashed, green!60!black, ->] ($(CL.south west)+(-0.2,-0.2)$) to[out=220, in=-40]
    node[below, pos=0.5] {\itshape naturaleza del movimiento}
    ($(VS.south east)+(0.2,-0.2)$);
\end{tikzpicture}
\end{center}

\begin{center}
\footnotesize\textbf{Figura~\ref{fig:flujo_control}:} Retroalimentación conceptual entre la ley de control y la generación de referencias.
\end{center}


%%%%%%%%%%%%%%%%%%%%%%%%%%%%%%%%%%%%%%%%%%%%%%%%%%%%%%%%%%

%%%%%%%%%%%%%%%%%%%%%%%%%%%%%%%%%%%%%%%%%%%%%%%%%%%%%%%%%%

\section{Consenso}
\subsection{Construcción del Error de Formación}

Cada agente construye su error de formación, el cual es utilizado por la ley de control para calcular la velocidad de referencia.

\begin{center}
\begin{tikzpicture}[
  font=\fontfamily{ptm}\selectfont\scriptsize,
  >={Stealth[length=2mm]},
  block/.style={draw, minimum width=1.3cm, minimum height=1.3cm, align=center, inner sep=2pt},
]
  \node[block] (VS) {Estructura\\Virtual};
  \node[block, right=1.1cm of VS] (C)  {Consenso};
  \node[block, right=1.1cm of C]  (CL) {Ley de\\Control};
  \node[block, right=1.1cm of CL] (F)  {PVA};

  \node[align=center, above=0.7cm of C] (NB) {Vecindad};
  \node[align=center, above=0.7cm of F] (LI) {LiDAR};

  \draw[->] (VS) -- node[above]{$x_v$} node[below]{$[x,y]$} (C);
  \draw[->] (C)  -- node[above]{$z$} node[below]{$[a,\alpha]$} (CL);
  \draw[->] (CL) -- node[above]{$u_{ref}$} node[below]{$[v,\omega]$} (F);
  \draw[->] (NB) -- (C);
  \draw[->] (LI) -- (F);

  \coordinate (uEnd) at ($(F.east)+(0.8,0)$);
  \draw[->] (F.east) -- node[above]{$u^*$} node[below]{$[v,\omega]$} (uEnd);

  \draw[dashed, orange] ($(C.north west)+(-0.2,0.2)$) rectangle ($(C.south east)+(0.2,-0.2)$);
  \draw[dashed, orange] ($(CL.north west)+(-0.2,0.2)$) rectangle ($(CL.south east)+(0.2,-0.2)$);

  \coordinate (CrectSouth) at ($(C.south)+(0,-0.2)$);
  \coordinate (CLrectSouth) at ($(CL.south)+(0,-0.2)$);
  \coordinate (arcMid) at ($(CrectSouth)!0.5!(CLrectSouth)+(0,-0.85)$);

  \draw[dashed, orange, ->] (CrectSouth) to[bend right=45]
    (CLrectSouth);

  \node[orange] at (arcMid) {\itshape error de formación};
\end{tikzpicture}
\end{center}

\begin{center}
\footnotesize\textbf{Figura~\ref{fig:flujo_control}:} Del error de consenso a la ley de control.
\end{center}

%%%%%%%%%%%%%%%%%%%%%%%%%%%%%%%%%%%%%%%%%%%%%%%%%%%%%%%%%%




\section{Consenso}
\subsection{Error de Formación con Estructura Virtual}

Se define un desplazamiento $\mathbf{r}_i$ para cada agente $i$, correspondiente a su posición deseada dentro de la estructura virtual:
\begin{equation*}
    \mathbf{e}_i = -\sum_{j \in \mathcal{N}_i} \big[(\mathbf{x}_i - \mathbf{r}_i) - (\mathbf{x}_j - \mathbf{r}_j)\big] - \beta\, b_i^0\big[(\mathbf{x}_i - \mathbf{r}_i) - \mathbf{x}_0\big],
\end{equation*}

El error de formación $\mathbf{e}_i = [e_x, e_y]$ se convierte a coordenadas polares, entrada de la ley de control:
\begin{equation*}
    a = \sqrt{e_x^2 + e_y^2}, \qquad \alpha = \operatorname{atan2}(e_y, e_x) - \theta_i.
\end{equation*}


%%%%%%%%%%%%%%%%%%%%%%%%%%%%%%%%%%%%%%%%%%%%%%%%%%%%




%%%%%%%%%%%%%%%%%%%%%%%%%%%%%%%%%%%%%%%%%%%%%%%%%%%%%%%%%%

\section{Ley de Control}
\subsection{Generación y Limitación de Señales de Referencia}

La ley de control genera las señales de referencia, las cuales son limitadas por el PVA en función de las lecturas de los obstáculos.

\begin{center}
\begin{tikzpicture}[
  font=\fontfamily{ptm}\selectfont\scriptsize,
  >={Stealth[length=2mm]},
  block/.style={draw, minimum width=1.3cm, minimum height=1.3cm, align=center, inner sep=2pt},
]
  \node[block] (VS) {Estructura\\Virtual};
  \node[block, right=1.1cm of VS] (C)  {Consenso};
  \node[block, right=1.1cm of C]  (CL) {Ley de\\Control};
  \node[block, right=1.1cm of CL] (F)  {PVA};

  \node[align=center, above=0.7cm of C] (NB) {Vecindad};
  \node[align=center, above=0.7cm of F] (LI) {LiDAR};

  \draw[->] (VS) -- node[above]{$x_v$} node[below]{$[x,y]$} (C);
  \draw[->] (C)  -- node[above]{$z$} node[below]{$[a,\alpha]$} (CL);
  \draw[->] (CL) -- node[above]{$u_{ref}$} node[below]{$[v,\omega]$} (F);
  \draw[->] (NB) -- (C);
  \draw[->] (LI) -- (F);

  \coordinate (uEnd) at ($(F.east)+(0.8,0)$);
  \draw[->] (F.east) -- node[above]{$u^*$} node[below]{$[v,\omega]$} (uEnd);

  \draw[dashed, violet] ($(CL.north west)+(-0.2,0.2)$) rectangle ($(CL.south east)+(0.2,-0.2)$);
  \draw[dashed, violet] ($(F.north west)+(-0.2,0.2)$) rectangle ($(F.south east)+(0.2,-0.2)$);

  \coordinate (CLrectSouth) at ($(CL.south)+(0,-0.2)$);
  \coordinate (FrectSouth) at ($(F.south)+(0,-0.2)$);
  \coordinate (arcMid) at ($(CLrectSouth)!0.5!(FrectSouth)+(0,-0.85)$);

  \draw[dashed, violet, ->] (CLrectSouth) to[bend right=45]
    (FrectSouth);

  \node[violet] at (arcMid) {\itshape señal de referencia};
\end{tikzpicture}
\end{center}

\begin{center}
\footnotesize\textbf{Figura~\ref{fig:flujo_control}:} De la ley de control a la limitación por el PVA.
\end{center}